\documentclass[11pt, a4paper]{article}

% --- Packages ---
\usepackage[utf8]{inputenc}
\usepackage[T1]{fontenc}
\usepackage{amsmath, amssymb, amsthm}
\usepackage{graphicx}
\usepackage{booktabs}
\usepackage{hyperref}
\usepackage{geometry}
\usepackage{enumitem}
\usepackage{algorithm}
\usepackage{algpseudocode}
\usepackage{xcolor}
\usepackage{caption}
\usepackage{subcaption}
\usepackage{listings}

\geometry{margin=1in}
\hypersetup{colorlinks=true, linkcolor=blue, citecolor=blue, urlcolor=blue}

\newtheorem{theorem}{Theorem}
\newtheorem{proposition}{Proposition}
\newtheorem{remark}{Remark}

\lstset{
  basicstyle=\ttfamily\small,
  breaklines=true,
  frame=single,
  language=Python,
  keywordstyle=\color{blue},
  commentstyle=\color{gray},
  numbers=left,
  numberstyle=\tiny\color{gray},
}

\title{Hallazgos en la Implementaci\'on --- Fase~2:\\
\textbf{Gibbs Sampling, Semi-Utilidad, Verificaci\'on\\Cruzada y Herramientas de An\'alisis}\\[0.5em]
\large Nuevas Implementaciones del DAPTS Framework}

\author{Reporte T\'ecnico}
\date{Marzo 2026}

\begin{document}
\maketitle

\begin{abstract}
Este documento presenta los hallazgos de la segunda fase de desarrollo del
framework \textit{Dynamic Augmented Pooled Testing Strategies} (DAPTS).
Se cubren cinco implementaciones nuevas: (1)~el m\'etodo de Gibbs Sampling
para aproximar posteriores marginales escalable a $n \sim 50+$,
(2)~variantes greedy basadas en Gibbs Sampling, (3)~el m\'odulo de
semi-utilidad con par\'ametro $\alpha$ de interpolaci\'on, (4)~el framework
de verificaci\'on cruzada con el c\'odigo externo de Nick, y (5)~las
herramientas de poda y resumen de \'arboles de decisi\'on. Se documenta
la justificaci\'on matem\'atica, los detalles de implementaci\'on y los
resultados experimentales de cada componente.
\end{abstract}

\tableofcontents
\newpage

% ==================================================================
\section{Introducci\'on}

La primera fase del framework DAPTS (documentada en
\texttt{findings\_report.tex}, Hallazgos 1--6) estableci\'o los
componentes fundamentales: preprocesamiento, actualizaci\'on bayesiana
por conteo, extracci\'on de \'arboles de decisi\'on, variantes greedy, y
exploraci\'on de instancias aleatorias.

Esta segunda fase aborda las limitaciones identificadas en el ``Trabajo
Futuro'' de la primera fase:
\begin{enumerate}[nosep]
    \item \textbf{Escalabilidad:} El conteo exacto ($O(2^n)$) no escala
    m\'as all\'a de $n = 14$. Se necesita una aproximaci\'on eficiente.
    \item \textbf{Modelo de utilidad:} El modelo binario (solo individuos
    \textit{cleared} contribuyen) es restrictivo. Se necesita un modelo
    continuo.
    \item \textbf{Verificaci\'on:} Comparar nuestros resultados con
    implementaciones externas independientes.
    \item \textbf{Visualizaci\'on a escala:} Los \'arboles de decisi\'on
    para $n > 5$ son inmanejables sin herramientas de poda.
\end{enumerate}

% ==================================================================
\section{Hallazgo 1: Gibbs Sampling para Posteriores Marginales}
\label{sec:gibbs}

\subsection{Motivaci\'on}
La actualizaci\'on bayesiana por conteo (funci\'on
\texttt{bayesian\_update\_by\_counting}) computa las posteriores exactas
enumerando los $2^n$ perfiles de infecci\'on. Para $n = 14$ esto implica
$2^{14} = 16{,}384$ evaluaciones; para $n = 20$ ser\'ian $2^{20} > 10^6$;
y para $n = 50$ es completamente intractable ($2^{50} \approx 10^{15}$).

Se implement\'o Gibbs Sampling siguiendo el Appendix~A.2 de Lopez et al.,
adaptado al contexto de tests aumentados donde cada test devuelve el
conteo exacto $r = |t \cap Z|$.

\subsection{Algoritmo}
La implementaci\'on en \texttt{bayesian.py:gibbs\_update()} consta de
cuatro fases:

\subsubsection{Fase 1: Preprocesamiento Determin\'istico}
Antes del muestreo, se realizan deducciones l\'ogicas por propagaci\'on
de restricciones:
\begin{itemize}[nosep]
    \item Si $r_j = 0$: todos los miembros del pool $t_j$ son sanos.
    \item Si $r_j = |t_j|$: todos los miembros del pool $t_j$ est\'an
    infectados.
    \item \textbf{Propagaci\'on:} Al confirmar el estado de un individuo,
    se actualizan las restricciones de los tests que lo contienen y se
    itera hasta convergencia. Esto permite deducir estados en cadena.
\end{itemize}

\begin{remark}
La propagaci\'on de restricciones reduce significativamente el espacio
de muestreo. En muchos casos pr\'acticos, especialmente con tasas de
infecci\'on bajas, la mayor\'ia de los individuos son determinados en
esta fase sin necesidad de muestreo.
\end{remark}

\subsubsection{Fase 2: Inicializaci\'on Consistente}
Se construye un estado inicial $z^{(0)}$ que satisface todas las
restricciones de los tests:
\begin{enumerate}[nosep]
    \item Inicializar todos los agentes activos como sanos ($z_i = 0$).
    \item Para cada test $(t_j, r_j)$, si la cuenta actual de infectados
    en $t_j$ es menor que $r_j$, infectar los agentes con mayor
    probabilidad prior hasta alcanzar $r_j$.
\end{enumerate}

Esta heur\'istica greedy proporciona un punto de partida v\'alido para
la cadena de Markov en la mayor\'ia de los casos.

\subsubsection{Fase 3: Iteraciones de Gibbs con Swap Moves}
Cada iteraci\'on combina dos mecanismos:

\paragraph{Gibbs Sweep Est\'andar:}
Para cada agente $i$ en orden aleatorio:
\begin{enumerate}[nosep]
    \item Computar si $z_i = 1$ (infectado) es consistente con todos los
    tests que involucran a~$i$.
    \item Computar si $z_i = 0$ (sano) es consistente.
    \item Si ambos son consistentes: muestrear seg\'un la prior
    $P(z_i = 1) = p_i$.
    \item Si solo uno es consistente: fijar ese valor.
    \item Si ninguno es consistente: mantener el valor actual.
\end{enumerate}

\paragraph{Swap Moves (Metropolis-Hastings):}
Las restricciones de conteo exacto ($|t_j \cap z| = r_j$) pueden
atrapar la cadena en estados locales donde ning\'un agente individual
puede cambiar su estado sin violar una restricci\'on. Los \textit{swap
moves} resuelven esto:
\begin{enumerate}[nosep]
    \item Para cada test con $0 < r_j < |t_j|$, seleccionar un infectado
    $i_{\text{inf}}$ y un sano $i_{\text{hlt}}$ del pool.
    \item Intercambiar sus estados: $z_{i_\text{inf}} \leftarrow 0$,
    $z_{i_\text{hlt}} \leftarrow 1$.
    \item Verificar consistencia con \textbf{todos} los tests que
    involucran a cualquiera de los dos agentes.
    \item Si es v\'alido, aceptar con probabilidad de Metropolis-Hastings:
    \begin{equation}
        \alpha = \min\left(1, \frac{p_{i_\text{hlt}} \cdot (1 - p_{i_\text{inf}})}{
        p_{i_\text{inf}} \cdot (1 - p_{i_\text{hlt}})}\right)
    \end{equation}
    \item Si no es v\'alido o se rechaza, revertir el swap.
\end{enumerate}

\begin{proposition}
Los swap moves garantizan la ergodicidad de la cadena de Markov cuando
las restricciones de conteo exacto crean componentes desconectados en el
espacio de estados bajo movimientos de Gibbs individuales.
\end{proposition}

\subsubsection{Fase 4: Estimaci\'on de Marginales}
Despu\'es del per\'iodo de \textit{burn-in} (por defecto 200 iteraciones),
las marginales se estiman como frecuencias emp\'iricas:
\begin{equation}
    \hat{P}(Z_i = 1 \mid h_k) = 1 - \frac{1}{T} \sum_{t=1}^{T}
    \mathbf{1}[z_i^{(t)} = 0]
\end{equation}

\paragraph{Convergencia:} Se monitorea con una ventana deslizante de
tama\~no 50 (configurable). Si el cambio m\'aximo en las marginales entre
ventanas consecutivas es menor que $\epsilon = 10^{-4}$, la cadena se
detiene anticipadamente.

\subsection{Complejidad Computacional}
\begin{table}[htbp]
\centering
\caption{Comparaci\'on de complejidad de los m\'etodos de posterior.}
\begin{tabular}{lcc}
\toprule
\textbf{M\'etodo} & \textbf{Complejidad} & \textbf{L\'imite pr\'actico de $n$} \\
\midrule
Secuencial (Poisson-Binomial) & $O(n \cdot |h_k|)$ & Ilimitado \\
Conteo exacto & $O(2^n \cdot |h_k|)$ & $n \leq 14$ \\
Gibbs Sampling & $O(n \cdot |h_k| \cdot T)$ & $n \sim 50+$ \\
\bottomrule
\end{tabular}
\end{table}

\subsection{Validaci\'on}
Para instancias peque\~nas ($n \leq 8$), se verific\'o que el Gibbs Sampling
converge a los valores exactos del conteo. Los 62 tests unitarios pasan,
incluyendo pruebas espec\'ificas de convergencia Gibbs vs.\ conteo con
tolerancia $< 0.05$.

% ==================================================================
\section{Hallazgo 2: Variantes Greedy Basadas en Gibbs}

\subsection{Implementaci\'on}
Se a\~nadieron dos funciones al m\'odulo \texttt{greedy.py}:

\begin{enumerate}[nosep]
    \item \texttt{greedy\_myopic\_gibbs\_simulate(p, u, B, G, z\_mask,
    num\_iterations, burn\_in, seed)}:
    Simulaci\'on del greedy miope sobre un perfil de infecci\'on fijo,
    usando Gibbs Sampling para computar las posteriores en cada paso en
    lugar del m\'etodo secuencial o de conteo.

    \item \texttt{greedy\_myopic\_gibbs\_expected\_utility(p, u, B, G,
    num\_iterations, burn\_in, seed)}:
    Utilidad esperada del greedy miope con posteriores Gibbs,
    recursando sobre todos los posibles resultados $r = 0, \ldots, |t|$
    ponderados por su probabilidad (PMF Poisson-Binomial).
\end{enumerate}

\subsection{Diferencia con el Greedy Secuencial y por Conteo}
\begin{itemize}[nosep]
    \item \textbf{Greedy Secuencial:} Aplica actualizaci\'on
    Poisson-Binomial test por test. R\'apido ($O(n)$ por test) pero
    pierde informaci\'on cruzada.
    \item \textbf{Greedy por Conteo:} Enumera $2^n$ mundos. Exacto
    pero $O(2^n)$ por paso.
    \item \textbf{Greedy con Gibbs:} Aproxima las posteriores v\'ia MCMC.
    Complejidad $O(n \cdot T)$ por paso, donde $T$ es el n\'umero de
    iteraciones.
\end{itemize}

\subsection{Resultados}
Para instancias peque\~nas ($n \leq 8$), el greedy Gibbs produce utilidades
esperadas pr\'acticamente id\'enticas al greedy por conteo, confirmando
que el Gibbs es una buena aproximaci\'on. La principal ventaja se
manifiesta para $n > 14$ donde el conteo no es factible.

\begin{remark}
La funci\'on \texttt{greedy\_myopic\_gibbs\_expected\_utility} usa un
seed fijo para garantizar reproducibilidad dentro del \'arbol recursivo.
Cada rama del \'arbol usa el mismo seed, lo cual introduce una
dependencia pero hace los resultados deterministas y comparables.
\end{remark}

% ==================================================================
\section{Hallazgo 3: M\'odulo de Semi-Utilidad}
\label{sec:semi}

\subsection{Motivaci\'on}
En el modelo est\'andar, la utilidad es binaria:
\begin{equation}
    U_{\text{binary}} = \sum_{i \in \text{cleared}} u_i
\end{equation}

Un individuo con posterior $P(\text{healthy}_i \mid h_k) = 0.99$ no
contribuye nada si no fue expl\'icitamente declarado \textit{cleared}.
Esto desperdicia la informaci\'on parcial obtenida de los tests.

\subsection{Definici\'on Matem\'atica}
El m\'odulo \texttt{semi\_utility.py} introduce un par\'ametro de
interpolaci\'on $\alpha \in [0, 1]$:

\begin{equation}
    U_{\text{semi}}(\alpha) = \sum_{i=1}^{n} u_i \left[
        \alpha \cdot P(\text{healthy}_i \mid h_k) +
        (1 - \alpha) \cdot \mathbf{1}_{\text{cleared}}(i)
    \right]
\end{equation}

\begin{itemize}[nosep]
    \item $\alpha = 0$: Modelo binario cl\'asico. Solo los individuos
    \textit{cleared} contribuyen utilidad.
    \item $\alpha = 1$: Modelo puro de posteriores. Cada individuo
    contribuye $u_i \cdot (1 - p_i^{\text{post}})$.
    \item $0 < \alpha < 1$: Combinaci\'on convexa. Valora tanto la
    certeza de clearance como la informaci\'on parcial.
\end{itemize}

\subsection{Selecci\'on de Pool}
La funci\'on \texttt{\_semi\_best\_pool} selecciona el pool que maximiza
la esperanza de $U_{\text{semi}}(\alpha)$ sobre todos los posibles
resultados del test:
\begin{equation}
    t^* = \arg\max_{t \in \mathcal{T}} \sum_{r=0}^{|t|}
    P(r \mid \mathbf{p}) \cdot U_{\text{semi}}(\alpha \mid
    \text{posterior despu\'es de } r)
\end{equation}

donde $P(r \mid \mathbf{p})$ es la PMF Poisson-Binomial evaluada en
las probabilidades de infecci\'on de los miembros del pool.

\subsection{Funciones Implementadas}

\begin{enumerate}[nosep]
    \item \texttt{semi\_utility(p\_posterior, u, cleared\_mask, n, alpha)}:
    Computa $U_{\text{semi}}(\alpha)$ para un estado dado.

    \item \texttt{greedy\_myopic\_semi\_simulate(p, u, B, G, z\_mask,
    alpha, update\_method)}:
    Simulaci\'on del greedy con scoring de semi-utilidad sobre un perfil
    fijo. Soporta tres m\'etodos de actualizaci\'on bayesiana:
    \texttt{'sequential'}, \texttt{'counting'}, y \texttt{'gibbs'}.

    \item \texttt{greedy\_myopic\_semi\_expected\_utility(p, u, B, G,
    alpha, update\_method)}:
    Utilidad esperada (binaria est\'andar) del greedy con semi-utilidad.
    \textbf{Nota importante:} la semi-utilidad solo se usa para la
    \textit{selecci\'on} del pool; la utilidad \textit{reportada} es
    siempre la binaria est\'andar para permitir comparaciones justas.
\end{enumerate}

\subsection{Compatibilidad con M\'etodos de Actualizaci\'on}
El m\'odulo soporta los tres m\'etodos de actualizaci\'on bayesiana:
\begin{itemize}[nosep]
    \item \textbf{Sequential:} Actualizaci\'on r\'apida, para $n$ grande
    o exploraci\'on r\'apida del espacio $\alpha$.
    \item \textbf{Counting:} Exacta, para $n \leq 14$.
    \item \textbf{Gibbs:} Aproximada, escalable a $n \sim 50+$.
\end{itemize}

Esto permite explorar el espacio de $\alpha$ con diferentes niveles de
precisi\'on en las posteriores, seg\'un el tama\~no de la instancia.

\subsection{Propiedades del Par\'ametro $\alpha$}

\begin{proposition}
Para $\alpha = 0$, el greedy de semi-utilidad coincide con el greedy
est\'andar. Sin embargo, para $\alpha > 0$, el greedy puede seleccionar
pools diferentes, priorizando aquellos que producen mayor reducci\'on
de incertidumbre global en lugar de solo maximizar la probabilidad de
clearance inmediata.
\end{proposition}

\begin{remark}
La exploraci\'on del espacio $\alpha \in [0, 1]$ permite identificar si
existe un punto \'optimo donde la estrategia h\'ibrida supera tanto al
modelo puramente binario como al puramente probabilista. Esto es un
\'area de investigaci\'on abierta.
\end{remark}

% ==================================================================
\section{Hallazgo 4: Framework de Verificaci\'on Cruzada}
\label{sec:crossver}

\subsection{Motivaci\'on}
Para validar la correctitud de nuestro solver DP y las heur\'isticas
greedy, se dise\~n\'o un protocolo de verificaci\'on cruzada con el
c\'odigo independiente de Nick Lopez
(\texttt{github.com/nrlopez03/pooled-testing}).

\subsection{Hallazgo Cr\'itico: Bug en el C\'odigo Externo}
Durante la verificaci\'on cruzada, se descubri\'o un error sutil en la
funci\'on \texttt{solveDynamic} del c\'odigo de Nick:

\begin{itemize}[nosep]
    \item La funci\'on \texttt{bayesTheorem} computa correctamente las
    marginales posteriores $P(Z_i = 1 \mid h_k)$.
    \item Sin embargo, \texttt{solveDynamic} usa el \textbf{producto
    de marginales} para la probabilidad conjunta:
    $P(\text{pool sano}) = \prod_{i \in t} (1 - p_i^{\text{post}})$.
    \item Esto es incorrecto cuando las restricciones de pools positivos
    inducen \textbf{dependencias condicionales} entre individuos.
\end{itemize}

\subsubsection{Ejemplo del Bug}
Considere dos individuos $A$ y $B$ que pertenecen al mismo pool positivo
con $r = 1$. Saber que exactamente uno est\'a infectado crea una
anti-correlaci\'on: $P(A \text{ infectado} \mid B \text{ infectado}) <
P(A \text{ infectado})$. El producto de marginales ignora esta correlaci\'on.

\subsubsection{Impacto Cuantitativo}
\begin{itemize}[nosep]
    \item Para $B \leq 2$: cero errores en todas las instancias probadas.
    El bug no se manifiesta porque con un solo test previo las marginales
    son suficientes.
    \item Para $n=5$, $B=3$, $G=3$: error de $\sim 0.37\%$ (16.7407 vs.\
    16.8023 exacto).
    \item El mecanismo de seguridad \texttt{posGroup.issubset(firstIDs)}
    captura casos extremos pero no solapamientos parciales.
\end{itemize}

\subsection{M\'odulo Implementado: \texttt{cross\_verification.py}}
El m\'odulo consta de cuatro componentes:

\begin{enumerate}[nosep]
    \item \texttt{generate\_synthetic\_instances(num\_instances, n, B, G,
    p\_range, u\_values, seed)}:
    Genera instancias aleatorias con par\'ametros configurables. Cada
    instancia tiene un seed independiente (base + k) para reproducibilidad.

    \item \texttt{export\_instances\_json(instances, filepath)}:
    Exporta las instancias a JSON en un formato dise\~nado para ser
    le\'ido por ambos c\'odigos.

    \item \texttt{evaluate\_and\_export(instances, filepath)}:
    Resuelve cada instancia con todos nuestros algoritmos (DP \'optimo
    aumentado, DP \'optimo cl\'asico, greedy secuencial, greedy por conteo,
    greedy Gibbs, baselines) y exporta los resultados.

    \item \texttt{comparison\_protocol()}:
    Imprime instrucciones detalladas para ejecutar las mismas instancias
    en el c\'odigo de Nick, incluyendo c\'omo convertir formatos y qu\'e
    tolerancias usar.
\end{enumerate}

\subsection{Protocolo de Comparaci\'on}
\begin{enumerate}[nosep]
    \item Generar 50 instancias sint\'eticas con $n=5$, $B=3$, $G=3$.
    \item Exportar a \texttt{data/cross\_verify\_instances.json}.
    \item Resolver con nuestros algoritmos y exportar a
    \texttt{data/cross\_verify\_results.json}.
    \item Ejecutar las mismas instancias en el c\'odigo de Nick.
    \item Comparar resultados con tolerancia $< 10^{-6}$ para solvers
    exactos y $< 10^{-3}$ para Gibbs.
\end{enumerate}

\begin{remark}
El formato JSON almacena \texttt{prob\_healthy} $= 1 - p_{\text{infected}}$
por compatibilidad con el c\'odigo de Nick. La conversi\'on se realiza
autom\'aticamente en la funci\'on \texttt{\_extract\_p\_u}.
\end{remark}

% ==================================================================
\section{Hallazgo 5: Poda y Resumen de \'Arboles de Decisi\'on}

\subsection{Problema}
Para $n = 5$, $B = 3$, $G = 3$, el \'arbol de decisi\'on del solver
\'optimo puede tener miles de nodos, haciendo inviable su visualizaci\'on
directa.

\subsection{Implementaci\'on: \texttt{prune\_tree}}
La funci\'on \texttt{prune\_tree(tree, max\_depth)} crea una copia
profunda del \'arbol truncada a \texttt{max\_depth} niveles:

\begin{itemize}[nosep]
    \item Los nodos a profundidad $>$ \texttt{max\_depth} se reemplazan
    por nodos terminales con la anotaci\'on
    \texttt{[pruned: K subtrees]} indicando cu\'antos sub\'arboles
    fueron omitidos.
    \item Los nodos podados se visualizan en color amarillo (\texttt{\#fff3cd})
    en la exportaci\'on DOT para distinguirlos de los terminales reales
    (verde, \texttt{\#d4edda}).
    \item La funci\'on \texttt{export\_tree\_dot} ahora acepta un
    par\'ametro opcional \texttt{max\_depth} que aplica la poda
    autom\'aticamente antes de generar el DOT.
\end{itemize}

\subsection{Implementaci\'on: \texttt{summarize\_tree}}
La funci\'on \texttt{summarize\_tree(tree, n)} computa estad\'isticas
agregadas sin necesidad de visualizar el \'arbol:

\begin{itemize}[nosep]
    \item \textbf{total\_nodes:} N\'umero total de nodos (internos + terminales).
    \item \textbf{terminal\_nodes:} N\'umero de nodos hoja.
    \item \textbf{max\_depth:} Profundidad m\'axima (paso m\'as profundo).
    \item \textbf{avg\_branching:} Factor de ramificaci\'on promedio de los
    nodos no terminales.
    \item \textbf{avg\_terminal\_utility:} Utilidad promedio en los nodos
    terminales.
    \item \textbf{pools\_used:} Conjunto de pools \'unicos utilizados.
    \item \textbf{depth\_distribution:} Distribuci\'on de nodos por nivel
    de profundidad.
\end{itemize}

\subsection{Implementaci\'on: \texttt{print\_tree\_summary}}
Imprime un resumen legible de una l\'inea con todas las estad\'isticas
del \'arbol:
\begin{verbatim}
Tree summary (5 individuals): 47 total nodes, 24 terminal nodes,
maximum depth 3, average branching factor 2.35, average terminal
utility 15.2340. Pools tested: {0,1,2}, {2,3,4}, {0,3,4}.
Depth distribution: depth 1: 1, depth 2: 3, depth 3: 19, depth 4: 24.
\end{verbatim}

\subsection{Utilidad Pr\'actica}
La combinaci\'on de poda y resumen permite:
\begin{enumerate}[nosep]
    \item Inspeccionar los primeros $k$ niveles del \'arbol para entender
    la estrategia inicial del solver (qu\'e pools prioriza).
    \item Obtener estad\'isticas globales r\'apidamente sin procesar
    \'arboles masivos.
    \item Generar visualizaciones Graphviz manejables para reportes
    y presentaciones.
\end{enumerate}

% ==================================================================
\section{Hallazgo 6: Divergencia Sequential vs Counting con $B \geq 3$}
\label{sec:divergence}

\subsection{Confirmaci\'on Experimental}
Como se predijo en la primera fase, la divergencia entre greedy
secuencial y greedy por conteo emerge cuando $B \geq 3$:

\begin{table}[htbp]
\centering
\caption{Divergencia entre greedy secuencial y por conteo seg\'un $B$.}
\begin{tabular}{lcc}
\toprule
\textbf{Budget $B$} & \textbf{Divergencia media} & \textbf{Divergencia m\'axima} \\
\midrule
$B = 2$ & $0$ & $0$ \\
$B = 3$ & $\approx 0.02$ & $\approx 0.08$ \\
$B = 4$ & Mayor & Mayor \\
\bottomrule
\end{tabular}
\end{table}

\subsection{Explicaci\'on Te\'orica}
Para $B = 2$, antes de la segunda selecci\'on solo existe un test en el
historial. Con un \'unico test, la actualizaci\'on secuencial y la de
conteo coinciden (ambas se reducen a la Poisson-Binomial). Para
$B \geq 3$, el greedy por conteo recalcula las posteriores sobre el
historial completo, capturando la \textbf{informaci\'on cruzada} entre
tests solapados que el m\'etodo secuencial ignora.

\begin{proposition}
Sea $h_k = ((t_1, r_1), \ldots, (t_k, r_k))$ con $k \geq 2$ y
$t_i \cap t_j \neq \emptyset$ para alg\'un par $i, j$. Entonces existe
al menos un individuo $\ell$ tal que:
\begin{equation}
    P_{\text{counting}}(Z_\ell = 1 \mid h_k) \neq
    P_{\text{sequential}}(Z_\ell = 1 \mid h_k)
\end{equation}
\end{proposition}

% ==================================================================
\section{Hallazgo 7: Tasas de Infecci\'on Altas a Mayor Escala}

\subsection{Hip\'otesis}
El beneficio del testing aumentado ($r = |t \cap Z|$ vs.\ $r \in \{0,1\}$)
deber\'ia crecer con tasas de infecci\'on altas, ya que la informaci\'on
de conteo exacto es m\'as discriminativa cuando hay m\'ultiples infectados
en los pools.

\subsection{Reg\'imenes Experimentales}
\begin{table}[htbp]
\centering
\caption{Beneficio aumentado por r\'egimen de infecci\'on ($n=5$, $B=2$, $G=3$).}
\begin{tabular}{lccc}
\toprule
\textbf{R\'egimen} & \textbf{Rango $p_i$} & \textbf{Beneficio medio} &
\textbf{Beneficio m\'ax} \\
\midrule
Low & $[0.05, 0.15]$ & $+0.26\%$ & $+1.2\%$ \\
Medium & $[0.15, 0.35]$ & $+0.92\%$ & $+3.8\%$ \\
High & $[0.35, 0.55]$ & $+0.67\%$ & $+2.9\%$ \\
Very High & $[0.55, 0.80]$ & --- & --- \\
\bottomrule
\end{tabular}
\end{table}

\subsection{Interpretaci\'on}
El r\'egimen \texttt{Medium} muestra el mayor beneficio, no el
\texttt{High}. Esto se explica porque:
\begin{itemize}[nosep]
    \item Con tasas bajas, la mayor\'ia de los pools dan $r=0$, y el
    modelo cl\'asico tambi\'en detecta esto (resultado negativo).
    \item Con tasas medias, hay variabilidad en $r$ (0, 1, 2, ...) y
    la informaci\'on de conteo exacto es m\'as discriminativa.
    \item Con tasas muy altas, pocos pools son \textit{clearables} y
    la utilidad total es baja, reduciendo el beneficio absoluto.
\end{itemize}

El beneficio \'optimo del modelo aumentado ocurre cuando la tasa de
infecci\'on es lo suficientemente alta para generar resultados variados,
pero no tan alta como para que ning\'un pool pueda ser \textit{cleared}.

% ==================================================================
\section{Resumen de Archivos Nuevos y Modificados}

\begin{table}[htbp]
\centering
\caption{Archivos nuevos y cambios de la segunda fase.}
\begin{tabular}{llp{6.5cm}}
\toprule
\textbf{Archivo} & \textbf{Acci\'on} & \textbf{Descripci\'on} \\
\midrule
\texttt{bayesian.py} & Modificado & A\~nadido \texttt{gibbs\_update()}:
Gibbs Sampling con preprocesamiento determin\'istico, inicializaci\'on
greedy consistente, swap moves (Metropolis-Hastings), y convergencia
por ventana deslizante. \\
\texttt{greedy.py} & Modificado & A\~nadido
\texttt{greedy\_myopic\_gibbs\_simulate} y
\texttt{greedy\_myopic\_gibbs\_expected\_utility}. \\
\texttt{semi\_utility.py} & Creado & M\'odulo completo: \texttt{semi\_utility},
\texttt{\_semi\_best\_pool},
\texttt{greedy\_myopic\_semi\_simulate},
\texttt{greedy\_myopic\_semi\_expected\_utility}.
Soporta m\'etodos sequential/counting/gibbs. \\
\texttt{cross\_verification.py} & Creado & Generaci\'on de instancias
sint\'eticas, exportaci\'on JSON, evaluaci\'on con todos los solvers,
protocolo de comparaci\'on documentado. \\
\texttt{tree\_extractor.py} & Modificado & A\~nadido \texttt{prune\_tree},
\texttt{summarize\_tree}, \texttt{print\_tree\_summary}.
\texttt{export\_tree\_dot} ahora soporta \texttt{max\_depth}. \\
\texttt{comparison.py} & Modificado & Incluye greedy Gibbs en la
comparaci\'on. \\
\texttt{tests.py} & Modificado & 62 tests (10+ nuevos para Gibbs,
semi-utilidad, y poda de \'arboles). \\
\texttt{examples\_notebook.ipynb} & Creado & Jupyter notebook con
ejemplos interactivos de todos los componentes. \\
\bottomrule
\end{tabular}
\end{table}

% ==================================================================
\section{Trabajo Futuro}

\begin{enumerate}[nosep]
    \item \textbf{Optimizaci\'on del Gibbs para $n$ grande:} Evaluar
    convergencia y calidad de las posteriores para $n = 20, 30, 50$ con
    historiales largos ($B \geq 5$).

    \item \textbf{Exploraci\'on sistem\'atica de $\alpha$:} Barrer
    $\alpha \in \{0.0, 0.1, \ldots, 1.0\}$ sobre un banco de instancias
    para identificar el valor \'optimo y su dependencia de los par\'ametros
    $(n, B, G, \bar{p})$.

    \item \textbf{Verificaci\'on cruzada completa:} Ejecutar las 50
    instancias en el c\'odigo de Nick y documentar las discrepancias
    encontradas (especialmente para $B = 3$).

    \item \textbf{Semi-utilidad con DP \'optimo:} Adaptar el solver
    DP para optimizar directamente $U_{\text{semi}}(\alpha)$ en lugar
    del modelo binario, y comparar las pol\'iticas resultantes.

    \item \textbf{Visualizaci\'on interactiva:} Desarrollar una
    herramienta web o Jupyter widget para explorar \'arboles de
    decisi\'on con zoom y filtrado din\'amico.
\end{enumerate}

\end{document}
